%************************************************
\chapter{Unconditional Flow Matching Model for X-Ray Images}\label{ch:exp1}
%************************************************

As discussed in the literature review, most recent techniques modify existing pre-trained generative models, often 
by adding extra components or introducing latent space disentanglement, to improve performance, interpretability, or compositionality  
without reducing the efficiency. However, these foundational models are generally trained on large, broad datasets rather than the 
medical data central to this thesis. To address this domain gap, I trained an unconditional flow matching model 
directly on a medical dataset, establishing a tailored pre-trained baseline for subsequent experiments. 
It is not intended to achieve SOTA performance, but rather to provide a baseline model to be used in the following experiments
to study disentanglement and compositionality in the medical domain. \\
\\
This study utilized the NIH-Chest-X-ray dataset (see Chapter \ref{ch:data} for more details) alongside the Representation Autoencoder 
(RAE) framework \cite{zheng2025diffusiontransformersrepresentationautoencoders}. Within this method, image latent representations are 
generated using a pre-trained visual model, specifically, a DINOv2-Base model configured with 4 register tokens 
\cite{oquab2024dinov2learningrobustvisual,darcet2024visiontransformersneedregisters}. The original RAE authors trained a decoder to 
construct an autoencoder from a large pre-trained visual model, yielding latents that offer richer semantic meaning and superior 
reconstruction quality compared to the SD-VAE latents typically used in state-of-the-art Stable Diffusion models. Subsequently, 
they trained a flow matching model on these RAE latents using a Diffusion Transformer (DiT), achieving good results but with a high 
computational cost because of the high-dimensional nature of the latents. To accommodate this problem, they also propose a modification 
of the DiT that incorporates architectural modifications inspired by Decoupled Diffusion Transformer (DDT) models, which 
prevent the need to scale the width across the entire backbone. When trained on ImageNet, this model achieved SOTA performance in both 
FID and IS metrics compared to other pixel and latent diffusion models. \\
\\
Because of the impressive generation quality of the model presented in the RAE paper and the quality and utility of the RAE latents, 
I decided to train a similar model on the NIH-Chest-X-ray dataset. I decided to train a DiT-XL (DINOv2-B) model insted of the improved 
DiT$^{DH}$-XL because I want to have a real bottleneck to perform further experiments. 

\begin{table}[ht]
    \centering
    \caption{Hyperparameter configuration used for training.}
    \label{tab:exp1_params}
    \renewcommand{\arraystretch}{1.25}
    \begin{tabular}{@{}llc@{}}
        \toprule
        \textbf{Category}       & \textbf{Hyperparameter}       & \textbf{Value}        \\
        \midrule
        \multirow{4}{*}{Optimization}
                                & Optimizer                     & AdamW                 \\
                                & Learning rate                 & $2 \times 10^{-4}$    \\
                                & Weight decay                  & $0.0$    \\
                                & Gradient clip norm            & $1.0$                 \\
        \midrule
        \multirow{3}{*}{Schedule}
                                & Warmup epochs                 & $40$             \\
                                & LR schedule                   & Linear      \\
                                & Final learning rate           & $2 \times 10^{-5}$    \\
        \midrule
        \multirow{4}{*}{Architecture}
                                & Input size                    & $16$                 \\
                                & Patch size                    & $1$                 \\
                                & In channels                   & $768$                 \\
                                & Hidden dimension              & $1152$                 \\
                                & Depth                         & $28$                  \\
                                & Attention heads               & $16$                   \\
                                & Dropout                       & $0.1$                 \\
        \midrule
        \multirow{3}{*}{Training}
                                & Batch size                    & $128$                 \\
                                & Mixed precision               & \texttt{float32}     \\
                                & EMA decay                     & $0.995$              \\
                                & Training epochs               & $450$                 \\
                                & Training steps                & $222$                 \\
        \bottomrule
    \end{tabular}
\end{table}


\begin{table}[ht]
    \centering
    \caption{Performance metrics of the DiT-XL (DINOv2-B) trained on the NIH-Chest-X-ray dataset.}
    \label{tab:exp1_results}
    \begin{tabular}{@{}lcccccc@{}}
        \toprule
        \textbf{Model}      & \textbf{Epochs}       & \#\textbf{Params}      & \textbf{FID}$\downarrow$      & \textbf{IS}$\uparrow$     & \textbf{Precision}$\uparrow$      & \textbf{Recall}$\uparrow$     \\
        \midrule
        DiT-XL   & --                    & --                    & --                            & --                        & --                    & --                    \\
        DiT$^{DH}$-XL   & --                    & --                    & --                            & --                        & --                    & --                    \\
        \bottomrule
    \end{tabular}
\end{table}
 


%*****************************************
%*****************************************
%*****************************************
%*****************************************
%*****************************************